\documentclass[sigconf]{acmart}
% Cleaning
\usepackage{pifont}
\pagestyle{fancy} % removes running headers
\settopmatter{printacmref=false} % Removes citation information below abstract
\renewcommand\footnotetextcopyrightpermission[1]{} % removes footnote with conference information in first column
\pagestyle{plain} % removes running headers
\fancyhead{}
\settopmatter{printacmref=false, printfolios=false}
% Copyright
\setcopyright{none}
% \setcopyright{acmcopyright}
% \setcopyright{acmlicensed}
% \setcopyright{rightsretained}
% \setcopyright{usgov}
% \setcopyright{usgovmixed}
% \setcopyright{cagov}
% \setcopyright{cagovmixed}
\settopmatter{printacmref=false} % Removes citation information below abstract
\acmDOI{}


\setlength{\parskip}{0pt}
\setlength{\parsep}{0pt}
\setlength{\headsep}{0pt}
\setlength{\topskip}{0pt}
\setlength{\topmargin}{0pt}
\setlength{\topsep}{0pt}
\setlength{\partopsep}{0pt}
%\linespread{0.95}
\usepackage{mdwlist}
\usepackage{amsmath}
\usepackage{booktabs} % professional table rules (\toprule etc.)
\usepackage{xurl}     % allow URLs to break anywhere (avoids overfull lines)

\begin{document}
\title{Representations of ETF and Fund Allocation Information:\\
A Historical Review and Comparative Analysis}

% TODO: replace with the real author list (all team members;
% the graduate student goes last). Team members known from the
% group chat: Mengchen Li, Shahmeer Ather Kakwan, Taimoor.
\author{Author One}
\affiliation{
  \institution{University of Victoria}
}
\email{author1@example.com}
\author{Author Two}
\affiliation{
 \institution{University of Victoria}
}\email{author2@example.com}
\author{Author Three}
\affiliation{
 \institution{University of Victoria}
}\email{author3@example.com}

\section{A Historical Review:}

Fund portfolio allocation data is a core information asset in global
financial markets. Its recording and representation have evolved from
unstructured paper narratives to standardized structured data and semantic
models, driven by regulatory transparency mandates, market efficiency
demands, and computer technology.

\subsection{Brief History}

In the past 90 years, fund allocation disclosure has evolved from
unstructured narrative reports to standardized, machine-readable formats.
Under U.S. regulation, the Investment Company Act of 1940 required
registered funds to disclose portfolio holdings in periodic shareholder
reports and prospectuses [7], and from 2004 in quarterly Form N-Q filings
[8]; for decades these records were plain-text tables designed for human
reading. The first U.S. ETF (SPDR S\&P 500 ETF Trust, SPY, 1993)
introduced daily portfolio transparency [1], formalized for U.S. ETFs in
2019 by SEC Rule 6c-11 [12], while the SEC's EDGAR system (fully phased in
by 1996) moved filings online [6]. Regulators then added machine-readable
formats: XBRL-tagged prospectus summaries in 2009 [9] and, in 2016, Forms
N-PORT and N-CEN, which replaced text-based reports with XML filings ---
N-PORT reports monthly position-level holdings data filed quarterly, and
N-CEN is an annual fund census [10, 11]. In parallel, the Financial
Industry Business Ontology (FIBO) standardized fund and security semantics
[2, 3], and Browne et al.\ demonstrated FIBO-based reporting of dated
fund-level information and allocation summaries; commercial databases such
as Morningstar aggregate these disclosures for cross-fund research [4, 5].
Because our sample ETFs are Canadian, the U.S. regulatory examples serve
as background: they establish the disclosure patterns --- dated,
structured, provider-published allocation percentages --- that motivate
our model.

\subsection{Summary}

Table 1 compares the representations most relevant to our model.
Throughout, we keep two concepts distinct: Holding, a customer's
investment in an ETF at a portfolio observation date, and AllocationEntry,
a disclosed percentage (e.g., ``technology: 30\%'') within a dated
AllocationSnapshot published by the ETF provider.

\begin{table*}[!t]
\caption{Comparison of fund allocation representation paradigms}
\label{tab:paradigms}
\setlength{\tabcolsep}{4pt}
\begin{tabular}{p{0.14\columnwidth}p{0.18\columnwidth}p{0.18\columnwidth}p{0.24\columnwidth}p{0.20\columnwidth}}
\toprule
Representation & What It Captures & What It Leaves Out & What We Can Learn for Our Design & Related ERD Entities\\
\midrule
Regulatory narrative filings (U.S.; N-SAR / N-Q [7, 8])
& Audited holdings at reporting dates; the fund--security relationship
& Unstructured; manual extraction; holdings not linked to allocation categories
& Keep a disclosure separate from its entries; our model records disclosed percentages, not full holdings
& ETF; AllocationSnapshot (dated disclosure)\\
Daily ETF website disclosures (Rule 6c-11 [12, 1])
& Provider-published daily basket: securities, weights, publication date
& No uniform schema; no category or classification links; source versions not tracked
& Model each disclosure as AllocationSnapshot (ETF + date + dimension); preserve the source in SourceDocument
& ETFProvider, AllocationSnapshot, AllocationEntry, AllocationDimension, SourceDocument\\
Structured regulatory filings (U.S.; N-PORT / N-CEN [10, 11])
& N-PORT: standardized monthly position-level holdings; N-CEN: annual fund census
& Regulatory schedules and a 60-day lag; not daily provider disclosures
& Structured, dated disclosure is feasible; we model provider-published percentages, not filing schedules
& ETF; AllocationSnapshot (structured, dated)\\
FIBO semantic ontology (incl. Browne et al.\ [2, 3])
& Formal fund and security semantics; Browne et al.: dated fund information and allocation summaries
& Conceptual only; no instance-level allocation entries
& Reuse core semantics; add an explicit pattern linking ETF, date, entries, categories, classification context, source
& ETF, ETFIdentifier, ClassificationScheme, AllocationCategory\\
Commercial aggregated databases (e.g., Morningstar [4, 5])
& Normalized holdings and category data across funds
& Proprietary classification logic; limited per-entry source and date traceability
& Model ClassificationScheme as a separate entity and keep source links
& ETFProvider, ClassificationScheme, AllocationCategory, SourceDocument\\
\bottomrule
\end{tabular}
\end{table*}

\subsection{Summary: Implications for Our Model}

\begin{table}[h]
\centering
\begin{tabular}{p{0.95\columnwidth}}
\toprule
The comparison informs our ERD in three ways:\\
\midrule
(1) Holding (a customer's ETF weight at a holding date) and
AllocationEntry (a disclosed percentage within a dated AllocationSnapshot)
are modeled as distinct entities;\\
(2) each AllocationSnapshot links an ETF to an allocation date, dimension,
entries, and --- where the source provides it --- a classification scheme,
with SourceDocument preserving the captured source;\\
(3) existing representations cover parts of this pattern (e.g., Browne et
al.'s FIBO-based reporting of dated fund information), so our contribution
is an explicit integration pattern rather than a claim that prior work
lacks dates or weights.\\
\bottomrule
\end{tabular}
\end{table}

\section{References}

% The numbering below is Mason's own list [1]--[12] and is kept verbatim.
% The final team assembly may renumber or switch to \bibliography{bibliography}.

[1] State Street Global Advisors. SPDR S\&P 500 ETF Trust (SPY). Available
at \url{https://www.ssga.com/us/en/individual/capabilities/spdr-core-equity-etfs/spy-sp-500}

[2] EDM Council. FIBO --- Financial Industry Business Ontology. Available
at \url{https://edmcouncil.org/frameworks/industry-models/fibo/}

[3] Object Management Group. 2017. Financial Industry Business Ontology
(FIBO): Fund Business Concepts (FBC) Specification. Available at
\url{https://www.omg.org/spec/EDMC-FIBO/FBC/}

[4] Morningstar, Inc. Morningstar Category Classification Methodology.
Available at
\url{https://morningstardirect.morningstar.com/clientcomm/Morningstar_Category_Classification.pdf}

[5] Morningstar, Inc. Morningstar Data: Fund Holdings. Available at
\url{https://docs-analyticslab.morningstar.com/1.7.1/morningstar_data/holdings.html}

[6] U.S. Securities and Exchange Commission. Electronic Data Gathering,
Analysis, and Retrieval (EDGAR) System --- Regulatory Overview. Available
at \url{https://www.sec.gov/info/edgar/regoverview.htm}

[7] U.S. Securities and Exchange Commission. 1940. Investment Company Act
of 1940. Available at \url{https://www.sec.gov/about/about-securities-laws}

[8] U.S. Securities and Exchange Commission. 2004. Shareholder Reports and
Quarterly Portfolio Disclosure of Registered Management Investment
Companies (Final Rule; Form N-Q). Available at
\url{https://www.sec.gov/rules/2004/02/shareholder-reports-and-quarterly-portfolio-disclosure-registered-management}

[9] U.S. Securities and Exchange Commission. 2009. Interactive Data for
Mutual Fund Risk/Return Summary (Final Rule; XBRL tagging). Available at
\url{https://www.federalregister.gov/documents/2009/02/19/E9-3359/interactive-data-for-mutual-fund-riskreturn-summary}

[10] U.S. Securities and Exchange Commission. 2016. Investment Company
Reporting Modernization (Final Rule; Forms N-PORT and N-CEN). Available at
\url{https://www.federalregister.gov/documents/2016/11/18/2016-25349/investment-company-reporting-modernization}

[11] U.S. Securities and Exchange Commission. 2016. Investment Company
Reporting Modernization Rules --- Small Entity Compliance Guide. Available
at
\url{https://www.sec.gov/divisions/investment/guidance/secg-investment-company-reporting-modernization-rules}

[12] U.S. Securities and Exchange Commission. 2019. Exchange-Traded Funds
(Final Rule; Rule 6c-11). Available at
\url{https://www.federalregister.gov/documents/2019/10/24/2019-21250/exchange-traded-funds}

\end{document}
